% !TeX root = ../papers/nc-p3.tex
\section{階数三の直交束により平面因子と二次因子を除く}
\label{sec:low-degree-octic-spinor}
\begin{theorem}\label{thm:low-degree-octic-spinor}
四辺有限の標準有限Frobenius窓で $X$ は被約、全階数反対平方根の階数2以下の軌跡が有限とする。
反対八次式は整であるか、既約五次式と既約三次式の積である。
\end{theorem}
\begin{proof}[三次の核ベクトルから四切断の零点へ]
次数一または二の既約因子 $S$ には三重の印がない。定理\ref{thm:quartic-octic-splitting}の証明より
核線束は真変換上で $P$ に等しく、四つの周囲の三次式 $u$ で核を表示できる。
$S$ の滑らかな印でこれらの生成イデアルは極大イデアルである。例外直線への写像が次数1の閉埋込みなので
一次部が余接空間を張るためである。ほかの基点は印に限られる。以下 $S$ またはその指定する引戻しの
滑らかな有理曲面を $T$、超平面の引戻しを $H$、$u$ の基底スキームを $Z$、長さを $s$ とする。
\[
 0\longrightarrow M\longrightarrow\mathcal O_T^4\xrightarrow{u^t}\mathcal I_Z(3H)\longrightarrow0 .
\]
$M$ は階数三の束である（正則曲面上の有限余長イデアルの射影次元は高々1）。
対称形式 $N$ は $Z$ の外で $M^*(-2H)\simeq M$ を与える。$N$ の核が $u$、階数が3だからである。
逆射とともに有限集合を越えて延長でき、$E=M(H)$ は行列式自明な直交三束となる。
\[
                       c_2(E)=6H^2-s .
\]
これは $c_1(M)=-3H$, $c_2(M)=9H^2-s$ から従う。$T$ は射影平面または有理線織曲面で
Brauer群が零なので、$\mathrm{SO}_3=\mathrm{PGL}_2$ の枠束は二束 $V$ から来て $E\simeq\operatorname{End}_0 V$ である。
ここで使うBrauer群の消滅は射影束公式と代数閉体上の $\mathbb P^1,\mathbb P^2$ の計算からも従う。
従って $c_2(E)=4c_2(V)-c_1(V)^2$。$M\hookrightarrow\mathcal O^4$ より任意の直線または母線
$J\simeq\mathbb P^1$ で $V|_J$ の二次数の差は $H J$ 以下である。実際 $E(-H)|_J$ は一般点で単射する。

$P_V\to T$ は $\pi_*\mathcal O(\xi)=V$ の規約の射影束とし、
$D=2\xi+H-c_1(V)$ と置く。$M^*\simeq E(H)=\operatorname{Sym}^2 V(H-c_1(V))$ であるから、
四つの座標は $P_V$ 上の $D$ の四切断になる。$N$ はそのGram行列で、各ファイバーの等方円錐は二項二次式の判別式である。
四切断の共通零点は、四つの二項二次式が共通の根を持つ点である。これは四つの値の直交空間に入る等方直線と同じで、
階数3の点にはなく、$\mathcal I_Z$ の局所生成数が2の印ではちょうど一つある。
実際 $M\to\mathcal O^4$ の値の階数は $4-2=2$ で、その二平面の制限形式は $N$ と同じ階数1である。
生成数3の印なら値は非等方な直線（$N$ の階数が1）なので、直交平面上に異なる二つの等方直線がある。
特に共通零点は有限である。四切断は独立である。依存すれば固定した非零ベクトルを $N$ が $S$ 上で消し、
一般核が一定となるが $L^2=9\deg S-m(S)>0$ である。

一般の三切断の完全交叉 $Y$ は全共通零点を含み、射影束公式により
\[
              \ell(Y)=D^3=6H^2+2(c_1(V)^2-4c_2(V))=2s-6H^2 .
\]
一般三切断が空の場合も交点数は零である。後に $H^0(\mathcal I_Y(D))$ が三定義切断に尽きることを使う。
そのためには $H^1(-D)=H^2(-2D)=0$ で十分である（Koszul分解、負の相対次数の零次切断は零）。
相対双対性からこれらは $H^0(T,-H)$ と $H^1(T,E(-2H))$ にそれぞれ等しい。
\end{proof}
\begin{proof}[平面と滑らかな二次曲面]
平面では $s=n=|X\cap S|\le7$。$2n-6\ge n$ により $n=6,7$ に限られる。
$n=6$ なら $c_2(E)=0$ で $c_1(V)$ は偶数、直線上の次数差は高々1なので全ての直線で均等である。
捩りを除けば $V$ は自明になる。実際任意の一点を吹き上げその点を通る直線の射影束と見ると、
引戻しは各ファイバーで自明ゆえ底からの引戻しであり、例外切断上で自明なため底でも自明である。
よって $H^1(E(-2H))=0$。$\ell(Y)=n$ で四番目の切断も $Y$ 全体で消えるが独立性に矛盾する。

$n=7$ なら $c_2(E)=-1$、$c_1(V)=1$, $c_2(V)=0$ に正規化できる。
$V|_J$ は各直線で $(1,0)$ 型ゆえ $V=\mathcal O\oplus\mathcal O(1)$。
実際Riemann--Rochで $\chi(V(-1))=1$、直線上の次数と双対性より $H^2(V(-1))=0$。
非零切断は線上の次数から零因子を持たず、零点も $c_2(V(-1))=0$ により存在せず、得た線束拡張は分裂する。
$P_V$ は一点での $\mathbb P^3$ の吹上げで $\xi$ は吹下げの超平面、$D=2\xi$。
従って独立な四つの二次式が $\mathbb P^3$ で有限な七点以上を切ることになる（吹下げの中心が基点なら
例外面全体が基点になってしまう）。一般三式の完全交叉長は8で四番目はその上で非零。
その零化イデアルで取る残余の長さは高々1だが、完全交叉の双対化束は $\mathcal O(2)$ で、
周囲の二次式は完全跡で消える。非零の切断の残余長が1なら、その切断倍の完全跡が定数1を消せず矛盾する。

滑らかな二次曲面では $s=n\le13$ で $2n-12\ge n$。二次曲面の交叉形式は偶なので $c_2(E)=12-n$ は偶数、
従って $n=12$ である。$c_1(V)=(a,b)$ で $2ab\equiv0\pmod4$ のため一方が偶数。
その母線上で二次数の差は高々1なので、捩れば全て自明である。従って $V$ は一方の $\mathbb P^1$ 上の束の
引戻しとなり、他の母線上でも次数差は高々1である。$E$ は三本の線束の和で、各二次数は一方が0、他方が $0,\pm1$ のいずれか。
$-2H$ を加えると両次数が負なので $H^1(E(-2H))=0$。再び $\ell(Y)=n$ は独立性に反する。
\end{proof}
\begin{proof}[二次錐の頂点]
頂点が印でないとき $T=\mathbb F_2$ に解消し、$H=C_0+2f$ とする。頂点で $N$ は階数3なので
例外切断 $C_0$ 上に基点はなく、$s=n\le13$。交叉形式が偶で $2n-12\ge n$ より再び $n=12$。
$C_0$ 上の二次数は等しい（$HC_0=0$）。$c_1(V)=a C_0+b f$ と書けば $b-2a$ は偶数であり、
$-2a^2+2ab$ が4で割れるので $a$ も偶数である。母線上で均等な $V$ は捩りを除き底から来て、
$C_0$ 上でも均等なので $E$ は自明。$H^1(-2H)=0$ は線織射の相対双対性による
（$R^1\pi_*\mathcal O(-2C_0-4f)=\mathcal O_{\mathbb P^1}(-2)$）。
これで $\ell(Y)=n$ の矛盾を得る。

頂点が印なら、錐を $x_0x_2=x_1^2$ とし、二重の有限射
\[
           T=\mathbb P^2\longrightarrow S,\qquad
             [v:t:w]\longmapsto[v^2:vt:t^2:w^2]
\]
を使う。第四座標は頂点以外の印が分岐平面を避けるように選ぶ。$H=2h$ と書く。
頂点の逆像で三次の核座標が生成するイデアルは $(v,t)^2$ である。実際元の頂点の接二次曲線では
核座標の一次部が全平面を張るため、引戻しの二次部は $v^2,vt,t^2$ を張り、中山の補題を使える。
他の $2(n-1)$ 個の印の逆像は単純、従って $s=2n+1$。頂点の逆像で生成数は3であるから
$P_V$ 上の基点は少なくとも $2(n-1)+2=2n$ 個である。$4n-22\ge2n$ より $n\ge11$。
$c_2(E)=23-2n$ が $4c_2(V)-c_1(V)^2$ として奇数なら4を法として3なので $n$ は偶数であり、$n\le13$ から12となる。
よって $c_2(E)=-1$、直線上で奇数の次数差が高々2なのでちょうど1である。
平面七点の場合と同じ証明から $V=\mathcal O\oplus\mathcal O(1)$ に正規化できる。

今は $P_V$ が同じ吹上げでも $D=2\xi+h$ で、完全交叉 $Y$ の長さは26、相異なる基点は24以上である。
$E$ が分裂して $H^1(E(-2H))=0$ なので、四番目の切断 $g$ は $Y$ 上で非零。
$\operatorname{ann}(g)$ で定まる残余 $W$ の長さは1または2である。$\omega_Y=(K_{P_V}+3D)|_Y$ で
$K_{P_V}+2D=2\xi$。従って $2\xi$ の切断は $W$ に全射できない。実際 $g$ 倍して完全交叉の跡を取れば
残余の完全跡であり、周囲から来る標準切断はKoszulの最終接続写像で消える。
ところが $2\xi$ は吹下げの二次Veroneseであり、長さ2以下を分離できないのは
$W$ が吹下げの例外面内に長さ2で入る場合だけである。
例外面上で $D$ は $\mathcal O_{\mathbb P^2}(1)$ となる。四つの一次式は共通基点が有限なので少なくとも二次元を張り、
一般三式のこの面上の共通零スキームは長さ高々1である。$W\subset Y$ と矛盾する。
平面と全二次因子が排除され、定理\ref{thm:exclude-tetrahedral-quartic}と合わせて残る分割は $5+3$ のみとなる。
\end{proof}
